\documentclass{article}
%%%%% PROGRAMMING PACKAGES
\usepackage{expl3,xparse}     %LaTeX3
\usepackage{mfirstuc}    %first letter uppercase with \makefirstuc (used in todo notes)
\usepackage{settobox}     %get box dimensions with \settoboxwidth (used in fit environment) => pretty stupid package.. could be avoided

%%%%% FIGURES, BOXES, COLORS
\usepackage{graphics,graphicx,trimclip,adjustbox}     %figures and box management
\usepackage{color,xcolor}     %to use colors (used in the todonotes)
    
%%%%% ENUMERATE
\usepackage{enumitem}   %to use "enumerate"
    \setlist[enumerate,1]{label = $\bullet$}   %default label at bullet points
    
%%%%% TODO NOTES
\usepackage[textwidth=16.3mm,textsize=scriptsize]{todonotes}   %to put notes everywhere
\newcommand*{\newcommentator}[2][red!20]{%
    \expandafter\newcommand\csname #2\endcsname[1]{\todo[color=#1]{{\bf\makefirstuc{#2}:} ##1}}%
    \expandafter\newcommand\csname big#2\endcsname[1]{\todo[inline,color=#1]{{\bf\makefirstuc{#2}:} ##1}}}
\newcommentator{pablo}

%%%%% MY HYPERLINK AND REF/LABEL MACROS
\usepackage{mypreamble/linkref}

%%%%% MATH PACKAGES
\usepackage{amssymb,amsmath}     %math original package
    \numberwithin{equation}{section}
\usepackage{mathtools}     %improves AMSmath
\usepackage{stackrel}     %better stacking of math symbols
\usepackage{mleftright}     %defines \mleft and \mright
    \mleftright     %redefines \left as \mleft and \right as \mright
\usepackage{pgfcore}
\usepackage{tikz-cd}     %advanced and rigorous diagrams
    
%%%%% MY SPECIAL PACKAGES
\usepackage{mypreamble/mathenv}     %improves on AMSmath's equation
\usepackage{mypreamble/thmenv}     %replaces the AMSthm package

%%%%% MATH MACROS AND FONT PACKAGES
%\usepackage{fontspec}
\usepackage{mypreamble/mathcro}

%%%%% HYPERREF PACKAGE
\makeatletter
\pretocmd\refstepcounter{%
 \zref@setcurrent{counter}{#1}}{}{}
\makeatother
\usepackage[pdfencoding=auto]{hyperref}   %hyperlinks parts of the pdf
    \hypersetup{
        bookmarksopen=true,   %expends tree of bookmarks in Acrobat pdf reader
        bookmarksnumbered=true,  %show numbers of sections in bookmarks
        linktoc=all,   %both section and page numbers are links
        hidelinks   %links (color+border) are hidden
    }
\providecommand*{\definitionautorefname}{Definition}
\providecommand*{\notationautorefname}{Notation}
\providecommand*{\conventionautorefname}{Convention}
\providecommand*{\conjectureautorefname}{Conjecture}
\providecommand*{\assumptionautorefname}{Assumption}
\providecommand*{\remarkautorefname}{Remark}
\providecommand*{\observationautorefname}{Observation}
\providecommand*{\lemmaautorefname}{Lemma}
\providecommand*{\propositionautorefname}{Proposition}
\providecommand*{\corollaryautorefname}{Corollary}
\providecommand*{\exampleautorefname}{Example}
\providecommand*{\exerciseautorefname}{Exercise}
\providecommand*{\openquestionautorefname}{Open Question}

%%%%% OUTPUT LAYOUT PACKAGES
\usepackage{geometry}     %to set margins
\geometry{
    a4paper,
    total={170mm,257mm},
    left=20mm,
    top=20mm,
    marginparwidth=20mm,
}
%% Clearpage before sections but not before TOC
\let\oldsection\section
\renewcommand\section{\clearpage\oldsection}
\makeatletter
\renewcommand\tableofcontents{%
    \oldsection*{\contentsname\@mkboth{\MakeUppercase\contentsname}{\MakeUppercase\contentsname}}\@starttoc{toc}%
}
\makeatother    %ensures the TOC stays in front page

%%%%% TYPESET DEBUGGING
%\usepackage{showframe}     %to make margins appear
%\usepackage{lua-visual-debug}     %to make all boxes appear
%\overfullrule=1pt   %creates 1pt wide black bars next to any overfull hbox

%%%%% QUIVER SETUP
\usetikzlibrary{calc,decorations.pathmorphing}
% A TikZ style for curved arrows of a fixed height, due to AndréC.
\tikzset{curve/.style={settings={#1},to path={(\tikztostart)
    .. controls ($(\tikztostart)!\pv{pos}!(\tikztotarget)!\pv{height}!270:(\tikztotarget)$)
    and ($(\tikztostart)!1-\pv{pos}!(\tikztotarget)!\pv{height}!270:(\tikztotarget)$)
    .. (\tikztotarget)\tikztonodes}},
    settings/.code={\tikzset{quiver/.cd,#1}
        \def\pv##1{\pgfkeysvalueof{/tikz/quiver/##1}}},
    quiver/.cd,pos/.initial=0.35,height/.initial=0}

% TikZ arrowhead/tail styles.
\tikzset{tail reversed/.code={\pgfsetarrowsstart{tikzcd to}}}
\tikzset{2tail/.code={\pgfsetarrowsstart{Implies[reversed]}}}
\tikzset{2tail reversed/.code={\pgfsetarrowsstart{Implies}}}
% TikZ arrow styles.
\tikzset{no body/.style={/tikz/dash pattern=on 0 off 1mm}}
\usepackage{mathdots}
\usepackage{bookmark}

%%%%% TITLE PAGE
\title{Higher Morita Categories}
\author{Pablo Bustillo Vazquez}

\begin{document}

\maketitle
\tableofcontents
\newpage
\section{Goal and Ideas}

For a given enrichment context $\mathcal{V}\in \EAlg{n}(\PrL)$, we have a sequence of symmetric-monoidal functors

\begin{equation}
    \EAlg{n} \longto \Mor_n^\alpha \longto \Mor_n^\beta \longto \PrL_n \longto \Cat{n}^{\Colim}
\end{equation}
The first category is $(\infty,1)$-$\mathcal{V}$-enriched, the others $(\infty,n+1)$-$\mathcal{V}$-enriched. The question of finding $k$-dualizable objects in these should be:
\begin{enumerate}
    \item (essentially) Kozsul duality in the first one (maybe more clear with augmented algebras?),
    \item factorization homology in the Morita ones then smoothness/properness
    \item higher condensation of previous ones
    \item same
\end{enumerate}

Our goals are:

\begin{enumerate}
    \item Construct a most general higher morita category by successive enrchiment. Requirement: stable under factorization homology.

    \item Embed such Morita category into nPr such that it controls all dualizable objects/cells.

    \item Duals up to $n$ = adjoints, then duals at $n+1$ = finiteness, and duals at $> n+1$ = invertibility. Extend characterization to 1 $\beta$-level and recover \textbf{fusion} categories.

    \item Link with Kozsul duality? Which objects are part of a Kozsul pair?

    \item Link to K-theory?

    \item Compute factorization homology of Mor (which is the natural domain for.. factorization homology). This corresponds to finding its full pasting calculus.

    \item Compute Brauer spectrum for several enrichment context.
\end{enumerate}

\section{Notations and Assumptions}

We assume a working theory of $(\infty,n)$-categories, for example the iterated enrichment model, in which we have:
\begin{enumerate}
    \item for each $n\geq 0$, an $(\infty, 1)$-category $\Cat{n}$,
    \item for each $n\geq 0$ and each $\mathcal{C} \in \Cat{n}$, a hom-functor
    \item $\mathcal{E}^k$-monoidal $(\infty,n)$-categories $\A$ have a $k$-delooping $\mathbb{B}^k\A$ which is an $(\infty,n+k)$-category,
    \item the $(\infty, 1)$
\end{enumerate}

\section{Dualizable Objects}

\subsection{Adjunctions}

\[
\begin{tikzcd}
    \cdots \arrow[r] & C_{n+1} \arrow[r, "d_{n+1}"] & C_n \arrow[r, "d_n"] & C_{n-1} \arrow[r] & \cdots \arrow[r] & C_1 \arrow[r, "d_1"] & C_0 \arrow[r] & 0
\end{tikzcd}
\]

\begin{definition}[Adjunction]
    Let $\mathcal{C}$ be an $(\infty,2)$-category. 
    A morphism $f\in a \to b \in \mathcal{C}$ is a left (resp. right) adjoint if it is a left (resp. right) adjoint in the bicategory $\tau_{\leq 2} \mathcal{C}$.
    We then write $f \dashv g$ to denote that $f$ is a left adjoint of $g$.
\end{definition}

We have the following very powerful result by Riehl and Verity, essentially estabilishing that the higher homotopical uniqueness of the adjunction data:

\begin{theorem}[Unicity of Adjoints]
    There is a $2$-category $\Adj$ together with the data of two morphisms $L, R \in \Adj^\to$ such that, for any $\mathcal{C}$ $(\infty,2)$-category, the precompositions
    % https://q.uiver.app/#q=WzAsMixbMCwwLCJbXFxBZGosXFxtYXRoY2Fse0N9XSJdLFsyLDAsIlxcbWF0aGNhbHtDfV5cXHRvIl0sWzAsMSwiTF4qIiwwLHsiY3VydmUiOi0xLCJzdHlsZSI6eyJ0YWlsIjp7Im5hbWUiOiJob29rIiwic2lkZSI6InRvcCJ9fX1dLFswLDEsIlJeKiIsMix7ImN1cnZlIjoxLCJzdHlsZSI6eyJ0YWlsIjp7Im5hbWUiOiJob29rIiwic2lkZSI6ImJvdHRvbSJ9fX1dXQ==
    \[\begin{tikzcd}
        {[\Adj,\mathcal{C}]} && {\mathcal{C}^\to}
        \arrow["{L^*}", curve={height=-6pt}, hook, from=1-1, to=1-3]
        \arrow["{R^*}"', curve={height=6pt}, hook', from=1-1, to=1-3]
    \end{tikzcd}\]
    induce \textbf{monomorphisms} of spaces selecting the connected components of left (resp. right) adjoint morphisms.
\end{theorem}

\textbf{Geometrically}, one can describe explicitly the $2$-category $\Adj$ geometrically as follows:
\begin{enumerate}
    \item \textbf{Objects:} the positively and negatively oriented point $\{\pm\}$,
    \item \textbf{Morphisms:} finite $1$-stratification of the real line where orientation of each segment alternates (we get sink and escape points in between them),
    \item \textbf{$2$-cells:} \textit{(up to isotopy)} finite $2$-stratification of the real plane where the oriented tangent vectors of the $1$-stratas is not allowed to be $(-1,0)$,
\end{enumerate}
\textbf{Combinatorially}, one can also describe explicitly the hom-categories as:
\[\begin{cases}[rCl]
    \Adj(+,+) &=& \Delta_+ = \textnormal{Free}^{\mathcal{E}^1}_{\mathcal{E}^1}(\{RL\})\\
    \Adj(+,-) &=& \Delta_\infty \\
    \Adj(-,-) &=& \Delta_+^\op = \textnormal{Free}^{\textnormal{co-}\mathcal{E}^1}_{\textnormal{co-}\mathcal{E}^1}(\{LR\})\\
    \Adj(-,+) &=& \Delta_{-\infty}
\end{cases}\]
where $\Delta_\infty$ (resp. $\Delta_{-\infty}$) is the category of finite non-empty linear orders with top (resp. bottom) element preserving morphisms.

Similarly, one can talk about $k$-morphisms in an $(\infty,n)$-category being adjoints.
Explicitly, if we let $\Adj_1 = \Adj$, one can build by induction the following ``walking adjunct pair of $k$-morphisms'' as the representing object in $\Cat{k+1}$ of the functor sending $\mathcal{C}\in\Cat{k+1}$ to the fibration of spaces classifying
% https://q.uiver.app/#q=WzAsNCxbMCwxLCJcXHRhdV8wXFxtYXRoY2Fse0N9XFx0aW1lcyBcXHRhdV8wXFxtYXRoY2Fse0N9Il0sWzEsMSwiXFxDYXR7a30iXSxbMCwwLCJcXENhdHtrKzF9KFxcQWRqX2ssIFxcbWF0aGNhbHtDfSkiXSxbNCwxLCJcXFNwYSJdLFswLDEsIlxcbWF0aGNhbHtDfSgtLC0pIiwyXSxbMiwwLCIiLDIseyJzdHlsZSI6eyJib2R5Ijp7Im5hbWUiOiJkYXNoZWQifX19XSxbMSwzLCJcXENhdHtrfShcXEFkal97ay0xfSwtKSIsMl1d
\[\begin{tikzcd}
	{\Cat{k+1}(\Adj_k, \mathcal{C})} \\
	{\tau_0\mathcal{C}\times \tau_0\mathcal{C}} & {\Cat{k}} &&& \Spa
	\arrow[dashed, from=1-1, to=2-1]
	\arrow["{\mathcal{C}(-,-)}"', from=2-1, to=2-2]
	\arrow["{\Cat{k}(\Adj_{k-1},-)}"', from=2-2, to=2-5]
\end{tikzcd}\]

We can also easily deduce unicity statements.
More explicitly, for $\Theta_k$ the $k$-globe, we have $L,R\in \Adj_k^{\Theta_k}$ $k$-cells, such that precomposition induces again monomorphisms
\[\begin{tikzcd}
    {[\Adj_k,\mathcal{C}]} && {\mathcal{C}^{\Theta_k}}
    \arrow["{L^*}", curve={height=-6pt}, hook, from=1-1, to=1-3]
    \arrow["{R^*}"', curve={height=6pt}, hook', from=1-1, to=1-3]
\end{tikzcd}\]

\subsection{Infinite Chain of Adjunctions}

\begin{definition}[Fully-Adjoint]
    Let $\mathcal{C}$ be an $(\infty,2)$-category and $f\in a\to b$ a morphism.
    We say that $f$ is part of an infinite chain of adjunctions if there exists a sequence of morphisms $f^{\vee n}$ such that:
    \[
        \cdots \dashv f^{\vee-n} \dashv \cdots \dashv f^{\vee-1} \dashv f \dashv f^{\vee1} \dashv \cdots \dashv f^{\vee n} \dashv \cdots
    \]
    We say that $f$ is infinitely-adjoint.
\end{definition}
If we let $\mathcal{C}^{\infty\textnormal{-adj}}$ be the following infinite fiber product,
% https://q.uiver.app/#q=WzAsMTAsWzQsMiwiW1xcQWRqLFxcbWF0aGNhbHtDfV0iXSxbMywzLCJcXG1hdGhjYWx7Q31eXFx0byJdLFsyLDIsIltcXEFkaixcXG1hdGhjYWx7Q31dIl0sWzYsMiwiW1xcQWRqLFxcbWF0aGNhbHtDfV0iXSxbMCwyLCJcXGNkb3RzIl0sWzgsMiwiXFxjZG90cyJdLFsxLDMsIlxcbWF0aGNhbHtDfV5cXHRvIl0sWzUsMywiXFxtYXRoY2Fse0N9XlxcdG8iXSxbNywzLCJcXG1hdGhjYWx7Q31eXFx0byJdLFs0LDAsIlxcbWF0aGNhbHtDfV57XFxpbmZ0eS1cXHRleHRub3JtYWx7YWRqfX0iXSxbMCwxLCJMIiwyLHsic3R5bGUiOnsidGFpbCI6eyJuYW1lIjoiaG9vayIsInNpZGUiOiJ0b3AifX19XSxbMiwxLCJSIiwwLHsic3R5bGUiOnsidGFpbCI6eyJuYW1lIjoiaG9vayIsInNpZGUiOiJib3R0b20ifX19XSxbMCw3LCJSIiwwLHsic3R5bGUiOnsidGFpbCI6eyJuYW1lIjoiaG9vayIsInNpZGUiOiJib3R0b20ifX19XSxbNCw2LCJSIiwwLHsic3R5bGUiOnsidGFpbCI6eyJuYW1lIjoiaG9vayIsInNpZGUiOiJib3R0b20ifX19XSxbMiw2LCJMIiwyLHsic3R5bGUiOnsidGFpbCI6eyJuYW1lIjoiaG9vayIsInNpZGUiOiJ0b3AifX19XSxbMyw3LCJMIiwyLHsic3R5bGUiOnsidGFpbCI6eyJuYW1lIjoiaG9vayIsInNpZGUiOiJ0b3AifX19XSxbMyw4LCJSIiwwLHsic3R5bGUiOnsidGFpbCI6eyJuYW1lIjoiaG9vayIsInNpZGUiOiJ0b3AifX19XSxbNSw4LCJMIiwyLHsic3R5bGUiOnsidGFpbCI6eyJuYW1lIjoiaG9vayIsInNpZGUiOiJ0b3AifX19XSxbOSwyLCIiLDAseyJjdXJ2ZSI6MSwic3R5bGUiOnsidGFpbCI6eyJuYW1lIjoiaG9vayIsInNpZGUiOiJib3R0b20ifX19XSxbOSwwLCIiLDAseyJzdHlsZSI6eyJ0YWlsIjp7Im5hbWUiOiJob29rIiwic2lkZSI6InRvcCJ9fX1dLFs5LDMsIiIsMCx7ImN1cnZlIjotMSwic3R5bGUiOnsidGFpbCI6eyJuYW1lIjoiaG9vayIsInNpZGUiOiJ0b3AifX19XSxbOSw0LCIiLDAseyJjdXJ2ZSI6Miwic3R5bGUiOnsidGFpbCI6eyJuYW1lIjoiaG9vayIsInNpZGUiOiJib3R0b20ifX19XSxbOSw1LCIiLDIseyJjdXJ2ZSI6LTIsInN0eWxlIjp7InRhaWwiOnsibmFtZSI6Imhvb2siLCJzaWRlIjoidG9wIn19fV1d
\[\begin{tikzcd}
	&&&& {\mathcal{C}^{\infty\textnormal{-adj}}} \\
	\\
	\cdots && {[\Adj,\mathcal{C}]} && {[\Adj,\mathcal{C}]} && {[\Adj,\mathcal{C}]} && \cdots \\
	& {\mathcal{C}^\to} && {\mathcal{C}^\to} && {\mathcal{C}^\to} && {\mathcal{C}^\to}
	\arrow[curve={height=12pt}, hook', from=1-5, to=3-1]
	\arrow[curve={height=6pt}, hook', from=1-5, to=3-3]
	\arrow[hook, from=1-5, to=3-5]
	\arrow[curve={height=-6pt}, hook, from=1-5, to=3-7]
	\arrow[curve={height=-12pt}, hook, from=1-5, to=3-9]
	\arrow["R", hook', from=3-1, to=4-2]
	\arrow["L"', hook, from=3-3, to=4-2]
	\arrow["R", hook', from=3-3, to=4-4]
	\arrow["L"', hook, from=3-5, to=4-4]
	\arrow["R", hook', from=3-5, to=4-6]
	\arrow["L"', hook, from=3-7, to=4-6]
	\arrow["R", hook, from=3-7, to=4-8]
	\arrow["L"', hook, from=3-9, to=4-8]
\end{tikzcd}\]
we can see that the projections down to any $\mathcal{C}^\to$ are monomorphisms of spaces selecting the connected components of infinitely-adjoint morphisms.

The functor $\mathcal{C} \mapsto \mathcal{C}^{\infty\textnormal{-adj}}$ should then be representable by some $2$-category that can be \textbf{geometrically} described as:
\begin{enumerate}
    \item \textbf{Objects:} the positively and negatively oriented point $\{\pm\}$,
    \item \textbf{Morphisms:} \pablo{to be described}
\end{enumerate}
This is however beyond the scope of this document.

Similarly to the previous subsection, one can describe the same notion for $k$-morphisms in an $(\infty,n)$-category of being infinitely-adjoint by taking the infinite fiber product
% https://q.uiver.app/#q=WzAsMTAsWzQsMiwiW1xcQWRqX2ssXFxtYXRoY2Fse0N9XSJdLFszLDMsIlxcbWF0aGNhbHtDfV57XFxUaGV0YV9rfSJdLFsyLDIsIltcXEFkal9rLFxcbWF0aGNhbHtDfV0iXSxbNiwyLCJbXFxBZGpfayxcXG1hdGhjYWx7Q31dIl0sWzAsMiwiXFxjZG90cyJdLFs4LDIsIlxcY2RvdHMiXSxbMSwzLCJcXG1hdGhjYWx7Q31ee1xcVGhldGFfa30iXSxbNSwzLCJcXG1hdGhjYWx7Q31ee1xcVGhldGFfa30iXSxbNywzLCJcXG1hdGhjYWx7Q31ee1xcVGhldGFfa30iXSxbNCwwLCJcXG1hdGhjYWx7Q31ee1xcaW5mdHlcXHRleHRub3JtYWx7LWFkai19a1xcdGV4dG5vcm1hbHstY2VsbH19Il0sWzAsMSwiTCIsMix7InN0eWxlIjp7InRhaWwiOnsibmFtZSI6Imhvb2siLCJzaWRlIjoidG9wIn19fV0sWzIsMSwiUiIsMCx7InN0eWxlIjp7InRhaWwiOnsibmFtZSI6Imhvb2siLCJzaWRlIjoiYm90dG9tIn19fV0sWzAsNywiUiIsMCx7InN0eWxlIjp7InRhaWwiOnsibmFtZSI6Imhvb2siLCJzaWRlIjoiYm90dG9tIn19fV0sWzQsNiwiUiIsMCx7InN0eWxlIjp7InRhaWwiOnsibmFtZSI6Imhvb2siLCJzaWRlIjoiYm90dG9tIn19fV0sWzIsNiwiTCIsMix7InN0eWxlIjp7InRhaWwiOnsibmFtZSI6Imhvb2siLCJzaWRlIjoidG9wIn19fV0sWzMsNywiTCIsMix7InN0eWxlIjp7InRhaWwiOnsibmFtZSI6Imhvb2siLCJzaWRlIjoidG9wIn19fV0sWzMsOCwiUiIsMCx7InN0eWxlIjp7InRhaWwiOnsibmFtZSI6Imhvb2siLCJzaWRlIjoidG9wIn19fV0sWzUsOCwiTCIsMix7InN0eWxlIjp7InRhaWwiOnsibmFtZSI6Imhvb2siLCJzaWRlIjoidG9wIn19fV0sWzksMiwiIiwwLHsiY3VydmUiOjEsInN0eWxlIjp7InRhaWwiOnsibmFtZSI6Imhvb2siLCJzaWRlIjoiYm90dG9tIn19fV0sWzksMCwiIiwwLHsic3R5bGUiOnsidGFpbCI6eyJuYW1lIjoiaG9vayIsInNpZGUiOiJ0b3AifX19XSxbOSwzLCIiLDAseyJjdXJ2ZSI6LTEsInN0eWxlIjp7InRhaWwiOnsibmFtZSI6Imhvb2siLCJzaWRlIjoidG9wIn19fV0sWzksNCwiIiwwLHsiY3VydmUiOjIsInN0eWxlIjp7InRhaWwiOnsibmFtZSI6Imhvb2siLCJzaWRlIjoiYm90dG9tIn19fV0sWzksNSwiIiwyLHsiY3VydmUiOi0yLCJzdHlsZSI6eyJ0YWlsIjp7Im5hbWUiOiJob29rIiwic2lkZSI6InRvcCJ9fX1dXQ==
\[\begin{tikzcd}
	&&&& {\mathcal{C}^{\infty\textnormal{-adj-}k\textnormal{-cell}}} \\
	\\
	\cdots && {[\Adj_k,\mathcal{C}]} && {[\Adj_k,\mathcal{C}]} && {[\Adj_k,\mathcal{C}]} && \cdots \\
	& {\mathcal{C}^{\Theta_k}} && {\mathcal{C}^{\Theta_k}} && {\mathcal{C}^{\Theta_k}} && {\mathcal{C}^{\Theta_k}}
	\arrow[curve={height=12pt}, hook', from=1-5, to=3-1]
	\arrow[curve={height=6pt}, hook', from=1-5, to=3-3]
	\arrow[hook, from=1-5, to=3-5]
	\arrow[curve={height=-6pt}, hook, from=1-5, to=3-7]
	\arrow[curve={height=-12pt}, hook, from=1-5, to=3-9]
	\arrow["R", hook', from=3-1, to=4-2]
	\arrow["L"', hook, from=3-3, to=4-2]
	\arrow["R", hook', from=3-3, to=4-4]
	\arrow["L"', hook, from=3-5, to=4-4]
	\arrow["R", hook', from=3-5, to=4-6]
	\arrow["L"', hook, from=3-7, to=4-6]
	\arrow["R", hook, from=3-7, to=4-8]
	\arrow["L"', hook, from=3-9, to=4-8]
\end{tikzcd}\]
where importantly, once again, all projections are monomorphisms.

\subsection{Iterated Adjunctions and Duals}

\begin{definition}[Fully-Adjoint]
    Let $\A$ be an $(\infty,n)$-category. By induction on $1\leq l$, we define the following. For $1 \leq k$, a $k$-morphism $f\in a_{k-1}\to b_{k-1}\in \ldots \in a_0 \to b_0$ is $l$-adjoint if:
\begin{enumerate}
    \item $f$ is part of an infinite chain of adjunctions in $(\A(a_0,b_0)\ldots)(a_{k-2}, b_{k-2})$ (interpreted as $\A$ when $k=1$)
    \[
        \cdots \dashv f^{\vee-n} \dashv \cdots \dashv f^{\vee-1} \dashv f \dashv f^{\vee1} \dashv \cdots \dashv f^{\vee n} \dashv \cdots
    \]
    with (co)units $\epsilon(f^{\vee k})$, $\eta(f^{\vee k})$ of $f^{\vee k} \dashv f^{\vee k+1}$ for $k \inset \mathbb{Z}$,
    \item (and if $l \leq 2$) for all $k \inset \mathbb{Z}$, $\epsilon(f^{\vee k})$ and $\eta(f^{\vee k})$ are $(l-1)$-adjoint.
\end{enumerate}
If $k+l = n$, we say that $f$ is fully-adjoint.
\end{definition}

\begin{remark}
    For $n + 1\geq k+l$, we recover the notion of invertibility.
\end{remark}

\begin{definition}[Fully-Dualizable Object] Let $\A$ be an $\mathcal{E}^1$-monoïdal $(\infty,n)$-category. For $1 \leq l$, an object $a \in \A$ is $l$-dualizable if it is $l$-adjoint as a $1$-morphism of $\mathbb{B}\A$. If $l=n$, we say that $a$ is fully-dualizable.
\end{definition}

\begin{definition}[Has Adjoints or Duals]\leavevmode
    \begin{enumerate}
        \item Let $\A$ be an $(\infty,n)$-category. We say that $\A$ \textbf{has adjoints} if all $k$-morphisms have a left and right adjoint for $1 \leq k \leq n-1$, or equivalently, if all $k$-morphisms are fully-adjoints.
        \item Let $\A$ be an $\mathcal{E}^1$-monoïdal $(\infty,n)$-category. We say that $\A$ \textbf{has duals} if $\mathbb{B}\A$ has adjoints.
    \end{enumerate}
\end{definition}

\begin{remark} If $\A$ is $\mathcal{E}^k$-monoïdal for $k \geq 2$, then an object (seen as a morphism of $\mathbb{B}\A$) has a left adjoint if and only if it has a right adjoint if and only if it is part of an infinite chain of adjunctions. We get identifications $a^{\vee 2n} \simeq a$, which can be made canonical for $\A$ symmetric-monoïdal.
\end{remark}

\begin{example}\label{rem:pointdual} $\Bordfr$ is symmetric-monoïdal and has duals. In particular, the point $\{+\}$ is a fully-dualizable object. 
\end{example}

For the next result, we need to briefly define a notation for the data involved in a $l$-adjunction:

\begin{definition}[Higher Adjunction Data]
    An adjunction data is a word $\omega$ using the symbols $\epsilon$, $\eta$, $-^{\vee n}$ (for any $n \inset \mathbb{Z}$) applied to $f$. The depth of an adjunction data is the number $\textnormal{depth}(\omega)$ of times either of the letters $\eta$ or $\epsilon$ were used.
\end{definition}

For example $\eta(\eta(\epsilon(f^{\vee -5}))^{\vee 12})$ is an adjunction data of depth $3$.

\begin{theorem}[Unicity]\label{thm:unicity}
    Let $\A$ be an $(\infty,n)$-category, $1 \leq k, l$, $f\in a_{k-1}\to b_{k-1}\in \ldots \in a_0 \to b_0$ a $l$-adjoint $k$-morphism $l$-adjoint and $\omega$ be a word 
    Let $\omega$ be any word using the symbols $\epsilon$, $\eta$, $-^{\vee n}$ (for any $n \inset \mathbb{Z}$) applied to $f$.
\end{theorem}

\subsection{Identifying Handlebodies}

\subsubsection{Just $\epsilon$}

Ok so we want $D^k \mono \mathbb{R}^k \times (0,1)$ with first coordinate canonical and second coordinate $\vec{x} \mapsto (1-r^2)/2$. 

\subsection{General}

Let $a + b \leq n$. We then have
\[
    \partial(D^a \times D^b) = S^a \times D^b \bigsqcup_{S^a \times S^b} D^a \times S^b
\]
Using a tubular neighbourhood $S^a \times S^b \times D^1 \mono D^a \times S^b$, we have
\[
    \partial(D^a \times D^b) = S^a \times D^b \bigsqcup_{S^a \times S^b} S^a \times S^b \times D^1 \bigsqcup_{S^a \times S^b} D^a \times S^b
\]

\subsection{Redundancy of Adjunction Data}

\begin{theorem}[Redundancy]\label{thm:redundancy}
    Let $\A$ be an $(\infty,n)$-category, $1 \leq k, l$, $f\in a_{k-1}\to b_{k-1}\in \ldots \in a_0 \to b_0$ a $l$-adjoint $k$-morphism.
    Then $f$ is $(l+1)$-adjoint if and only if......
\end{theorem}

\subsection{Finding Fully-Dualizable Objects}

\begin{remark} Recall that we have functors $\Bordfr_k \to \Bordfr_n$ given by extending the framing trivially. \textbf{Warning:} these functors are \textbf{not} truncations.
\end{remark}

\begin{lemma}[Constructing Full-Adjunctions]\label{lem:constr} Let $\A$ be a symmetric-monoïdal $(\infty,n+1)$-category together with a functor $F\in \Bordfr_n \to \A$ sending $\{+\}$ to some object $a$. Then $A$ is fully-dualizable if and only if, for every $n$-morphism
    \[
        f\in a_{n-1}\to b_{n-1}\in \ldots \in a_0 \to b_0\in \Bordfr_n
    \]
    $Ff$ is part of an infinite chain of adjunctions.
\end{lemma}
\begin{proof} For $n = 0$, $\Bordfr_0$ is the trivial category $\{\bullet\}$ and the statement is tautological. Now assume it true for $n \geq 0$. Clearly, $A$ (seen as a morphism in $\mathbb{B}\A$) is part of an infnite chain of adjunctions given by the image of $\Bordfr_1 \to \Bordfr_n \to \A$.
\end{proof}

\section{Alpha-Case: Objects}

Recall that $\Mor_n$ is an $(\infty,n+1)$-category.

\subsection{First Step: n-dualizability}

It's well-known that every $k$-morphisms of $\Mor_n$ is $(n-k)$-adjoints (but not necessarily fully-adjoint). Using factorization homology, we can make the $n$-dualizability data of an object more explicit: 

\begin{theorem}[Cite Claudia and John] Let $\mathcal{V}\in \EAlg{n}\left(\Cat{1}^{\textnormal{sifted}-\Colim}\right)$ . For $a \in \EAlg{n}(\mathcal{V})$, we have a functor
\[
    \int_{-} a\in \Bordfr_n \longto \Mor_n(\mathcal{V})
\]
from an $(\infty,n)$-category to an $(\infty,n+1)$-category.
\end{theorem}

By Lemma \ref{lem:constr}, we then have the following:

\subsection{Second Step: (n+1)-cells}

For $A$ to be $n+1$-dualizable, it suffices for the $n$-cells involved in its $n$-dualizable structure to have adjoints. A priori, the cobordism hypothesis's inductive proof relies on the claim that it's enough to check adjoints on some of these cells (the spheres). We can prove it formally in this context.

The strategy is the following:
\begin{enumerate}[label = (\roman*)]
    \item Any arbitrarily complicated $n$-cell $x\in a_{n-1}\to b_{n-1}\in \ldots \in a_0 \to b_0$ in the image of our functor, can be seen (after smoothing) as the data of a $n$-manifold with boundary $M$ and where the action is given by the canonical map of $\mathcal{E}^1$-algebra
    \[
        \int_{\partial M \times \mathbb{R}} A \longto \End\left(\int_M A\right)
    \]

    \item To find a dual, we need to proceed by handle decomposition and show that each handle has a dual. TODO: link handles to the half-spheres.
\end{enumerate}

\section{Alpha-Case: Higher Cells}

Lots of elements in the proof are based on the idea that we have an inductive definition of $\Mor_n$, which was proved formally by Rune Haugseng I believe (Thm $1.3$). For dualizability to be recognized as a finiteness condition, one simply need to find a type of $\infty$-categories left unchanged by this inductive process:

\begin{definition}[Good Categories] A $n$-good category is a stable presentable $\Disk_n$-$\infty$-category, \emph{i.e.} a $n$-disk algebra in the $\infty$-category of stable presentable categories with its usual symmetric tensor product. 
\end{definition}

\begin{lemma}[Categories of Bimodules are Good] Let $\mathcal{S}$ be a $n$-good category and let $a, b\in \Disk_n(\mathcal{S})$ be $n$-disk algebra objects. Then ${}_a\Bimod_b(\mathcal{S})$ is a $(n-1)$-good category.
\end{lemma}
\begin{proof} Stability and presentability should probably be proven in [HA]. The $(n-1)$-disk operation should be clear to define. 
\end{proof}


I need [HA.4.6.2.13] to show that "being a left adjoint" only depends on the codomain algebra. Then general manipulations of disk-operads show it only depends on the bar construction. Then if $\mathcal{S}$ is spectra, or maybe R-mod, we know it means finite module. In more general setup it'd be some absolute colimits maybe? (using the link between being absolute and having a dual).

We first need to identify which $n$-cells are dualizable. First, we need to define a higher notion of two-sided bar construction (arguably a ``stratified-$S^{n-1}$''-sided bar construction).

\begin{definition}[Higher Bar Construction]
Let us define a stratification of $\mathbb{R}^{n}$ as follows: for $0\leq k \leq n-1$, we have two strata of the form $\{0\}^k \times \mathbb{R}_{\epsilon} \times \mathbb{R}^{n-k-1}$ with respectively $\epsilon \inset \{>0, <0\}$ and they are respectively called $A_{k}$ and $B_{k}$, and we have the additional strata $\{0\}^{n}$ called $X$. We can this stratified space $\BarCon_n$. For any $x\in a_{n-1}\to b_{n-1}\in \ldots \in a_0 \to b_0$ in $\Mor_n(\mathcal{S})$, we can then define $\int_{\BarCon_n} x$ where the stratas are labelled in the obvious way.
\end{definition}

\begin{lemma}[Dualizability of the highest-level cells] A $n$-cell $x\in a_{n-1}\to b_{n-1}\in \ldots \in a_0 \to b_0$ is dualizable in $\Mor_n(\mathcal{S})$ if and only if $x$ is a finite $\int_{\BarCon_{n-1}} a_{n-1}$-left-module and a finite $\int_{\BarCon_{n-1}}b_{n-1}$-right-module in $\mathcal{S}$.
\end{lemma}
\begin{proof} By the iterative definition of $\Mor_n(\mathcal{S})$, we see that $x$ is dualizable whenever it is dualizable as a $1$-cell in $\Mor_1({}_{a_{n-2}}\Bimod_{b_{n-2}}(\ldots {}_{a_0}\Bimod_{b_0}(\mathcal{S})))$. By [Insert reference to HA], $x$ is then dualizable if and only if it is a finite $a_{n-1}$-left-module and a finite $b_{n-1}$-right-module in this iterated bimodule category. Equivalently, it has to be dualizable as a finite $\int_{\BarCon_{n-1}} a_{n-1}$-left-module and a finite $\int_{\BarCon_{n-1}}b_{n-1}$-right-module in $\mathcal{S}$ directly (the categories should be equivalent, would be interesting to prove carefully).
\end{proof}

\begin{theorem}[Dualizability in $\Mor_n$]\label{dual} A $k$-cell $x\in a_{k-1}\to b_{k-1}\in \ldots \in a_0 \to b_0$ is dualizable if and only if
\begin{enumerate}
    \item $x$ is a finite $\int_{\BarCon_{k-1}} a_{k-1}$-left-module and a finite $\int_{\BarCon_{k-1}}b_{k-1}$-right-module in $\mathcal{S}$,
    \item for all $0\leq m \leq n-k-1$, $x$ is a finite $\int_{X_m} x$-module where $X_m$ is the stratified space $\mathbb{R}^{m+1} \times \BarCon_{k-1}$ with an additional strata $S^m \subset \mathbb{R}^{m+1}\times \{0\}^{k-1}$ with $x$ supported on the $m$-sphere, either $a_{k-1}$ or $b_{k-1}$ supported on the inside or outside (one can further multiply the space by $\mathbb{R}^{n-m-k}$ which is at least $1$-dimensional, hence it gives an algebra).
\end{enumerate}
\end{theorem}
\begin{proof} Dualizability can be checked at the level of $1$-cells in $\Mor_{n-k+1}({}_{a_{k-2}}\Bimod_{b_{k-2}}(\ldots {}_{a_0}\Bimod_{b_0}(\mathcal{S})))$.
\end{proof}

What does the frame version look like? One has to use that $S^m\times \mathbb{R}$ is canonically $(m+1)$-framed.

\section{Defining Morita Categories}

For $\mathcal{V} \in \EAlg{a+1}(\Pr)$, we want
\[
    \Mor_{a+1}(\mathcal{V})(A, B) = \Mor_{(a-1)+1}([A^\op \times B, \mathcal{V}])
\]

For $\mathcal{V} \in \EAlg{2}(\Pr)$, we want
\[
    \Mor_{0+2}(\mathcal{V})(A, B) = \Mor_{0+1}((A^\op \times B)^{\textnormal{loc.}\Colim})
\]

For $\mathcal{W}$ a locally-cocomplete $(\infty,2)$-category, we have $\mathcal{W}\Pr = \left[\mathcal{W}, \Cat{1}^{\Colim}\right]^{\textnormal{compact}}$. If $\mathcal{W}$ is localy-presentable, we should have $\mathcal{W}\Pr = \left[\mathcal{W}, \Pr\right]$.

Locally, $\Mor_{0+2}(\mathcal{V}) \to 2\Pr^\mathcal{V}$ should be given by some
\[
    \Mor_{0+1}((A^\op \times B)^{\textnormal{loc.}\Colim}) \longrightarrow \left[A^\op \times B, \Pr\right]
\]

Assuming the recognition of dualizable objects still work, one should find that dualizable endomorphism in $2\Pr^\mathcal{V}$ are retracts of one from the morita context. For this to work, one need to recover the monoidal structure... 

\section{One Beta Level}

\begin{conjecture}
    Fully-dualizable $n$-monoidal rigid $(\infty,1)$-$\mathcal{V}$-categories are Morita-equivalent to a fully-dualizable $(n+1)$-monoidal algebra.
\end{conjecture}

Because $\cat{A}$ is a finite $\cat{A} \times \cat{A}^\op$-module, we have 

\begin{conjecture}
    Fully-dualizable $n$-monoidal rigid $(\infty,1)$-$\mathcal{V}$-categories are \textbf{pivotal}: there exists a natural maps $x \simeq x^{\vee 2}$. 
\end{conjecture}



Question: in what way is this compatible with the monoidal structure?

Idea of proof: one can show they are morita equivalen to an alpha-version $A$. The biggest such is the ``category of finite one-sided modules''. There everything is an absolute colimit of $A$. The dual should be the corresponding absolute limit of $A$. Maybe this helps? Alternate between limit and colimit.

\begin{conjecture}
    One can construct a trace for any $n$-monoidal rigid pivotal $(\infty,1)$-$\mathcal{V}$-category.
\end{conjecture}

How would that construction work? Pivotal means we have a forgetful enriched-functor $\int_{S^1} \mathcal{C} \to \mathcal{V}$. In a way that means: one can almost define cyclic tensor product but only on underlying $\mathcal{V}$-objects.

\begin{conjecture}
    Fully-dualizable $n$-monoidal rigid $(\infty,1)$-$\mathcal{V}$-categories are \textbf{spherical}: the two possible traces coincide. 
\end{conjecture}

What other higher notions of the sort are true?

These ideas don't seem to give a "smallness/finiteness" to the category. They're probably not sufficient to characterize fd. It might be a good first step in understanding them though.


\section{Azumaya over Algebraically Closed Field}

\begin{proposition}\label{p2'}
Let $k$ be algebraically closed and $B$ a $n$-Azumaya $k$-algebra. Then $B$ is Morita equivalent to the trivial $k$-algebra $k$.
\end{proposition}

\textit{Proof:} The Künneth Theorem can be seen as a symmetric monoïdal functor $kMod \to kVect^{gr}$ from an $(\infty,1)$-category to a $1$-category.


For this, we let $H^{*}(B)$ be the graded cohomology $k$-algebra of $B$. As we are working over a field, 
we have natural identifications of graded $k$-vector spaces. 
$$H^{*}(B\otimes_{k}B^{op})\simeq H^{*}(B)\otimes_{k}H^{*}(B)^{op}$$
$$H^{*}(\mathbb{R}\underline{Hom}_{k}(B,B))\simeq Hom^{*}_{k}(H^{*}(B),H^{*}(B)).$$
These identifications, and conditions $\mathbf{(Az-1)}$ and $\mathbf{(Az-2)}$ imply that 
$H^{*}(B)$ is itself an Azumaya $k$-algebra, and more precisely that $H^{*}(B)$, considered as 
a graded bi-$H^{*}(B)$-module, induces an equivalence between the categories
of graded modules
$$H^{*}(B)\otimes_{k}H^{*}(B)^{op}-Mod^{gr} \simeq  k-Mod^{gr}.$$
This implies that in particular that $H^{*}(B)$ is a graded projective 
$H^{*}(B)\otimes_{k}H^{*}(B)^{op}$-module, 
and therefore that $H^{*}(B)$ is a semi-simple graded $k$-algebra (i.e. that 
$H^{*}(B)-Mod^{gr}$ is semi-simple $k$-linear abelian category), and thus that any 
graded $H^{*}(B)$-module is graded projective. 

We have a graded functor 
$$\phi : D(B) \longrightarrow H^{*}(B)-Mod^{gr}$$
sending a $B$-dg-module $E$ to its cohomology $H^{*}(E)$. 

\begin{lemma}\label{l2}
The functor above $\phi$ is an equivalence of (graded) categories. 
\end{lemma}

\textit{Proof:} Let call an object $E\in D(B)$ \emph{graded free} if it is a (possibly infinite) 
sum of shifts of $B$. An object of $D(B)$ will then be called \emph{graded projective} if it is a
retract of a graded free object. It is clear by definition that the functor $\phi$ is fully faithful when 
restricted to graded free objects, and thus also when restricted to graded projective objects. In order to 
show that $\phi$ is fully faithful it is then enough to show that any object in $D(B)$ is graded projective. 

Let $E \in D(B)$, and let $H^{*}(E)$ be the corresponding graded $H^{*}(B)$-module, which is graded projective 
as we saw. Let $p$ be a projector on some graded free graded $H^{*}(B)$-module $P$ with image $H^{*}(E)$. 
If we write 
$$P=\bigoplus_{d}(H^{*}(B)[p]^{(I_{p})}),$$
then $p$ defines a projector $q$ in $D(B)$ on the graded free object
$$Q:=\bigoplus_{d}(B[p]^{(I_{p})}) \in D(B).$$
As $D(B)$ is Karoubian complete (see for instance \cite[Sous-lemme 3]{to4}), we now that this projector splits
$$q : \xymatrix{
Q \ar[r]^-{k} & E' \ar[r]^-{j} & Q,}$$
with $E' \in D(B)$. We claim that $E$ and $E'$ are isomorphic objects in $D(B)$. Indeed, by the choice of 
$p$, there exists a factorisation in $H^{*}(B)-Mod^{gr}$
$$p : \xymatrix{
P \ar[r]^-{r} & H^{*}(E) \ar[r]^-{i} & P.}$$
The images of $k$ and $j$ by $\phi$ provides another such factorisation
$$\phi(q)=p : \xymatrix{
P \ar[r]^-{\phi(k)} & H^{*}(E') \ar[r]^-{\phi(l)} & P.}$$
By the uniqueness of such factorisation, we have that 
the composite morphism 
$$\xymatrix{
H^{*}(E') \ar[r]^-{\phi(l)} & P \ar[r]^-{r} & H^{*}(E)}$$
is an isomorphism.
Now, from the definition of $\phi$, the morphism $r$ lifts uniquily to a morphism 
$r' : Q \longrightarrow E$
with $\phi(r')=r$. We thus get a morphism in $D(B)$
$$\xymatrix{E' \ar[r]^-{j} & Q \ar[r]^-{r'} & E,}$$
whose image by $\phi$ is $r\circ \phi(l)$, which we have seen to be an isomorphism. 
As the functor $\phi$ is conservative we have $E \simeq E'$ in $D(B)$. 

We thus have that all objects in $D(B)$ are graded projective, and therefore that 
$\phi$ is fully faithful. All the objects $H^{*}(B)[p]$ belong to the essential image of $\phi$, and 
therefore so do all graded free $H^{*}(B)$-module. As $\phi$ is fully faithful and as 
$D(B)$ is  Karoubian closed, we deduce that any 
graded projective $H^{*}(B)$-module also belongs to the essential image of $\phi$. But, as all
graded $H^{*}(B)$-module are graded projective we have that $\phi$ is essentially surjective and thus an 
equivalence of categories. 
\hfill $\Box$ \\

We use the previous lemma as follows. Let $E\in H^{*}(B)-Mod^{gr}$ be a simple object, which is
automatically finite dimensional over $k$, as  $H^{*}(B)$ is so.  
Because $k$ is algebraically closed we have 
$$Hom^{0}_{H^{*}(B)-Mod^{gr}}(E,E) \simeq k.$$
Moreover, for any $n\neq 0$, we have
$$Hom^{n}_{H^{*}(B)-Mod^{gr}}(E,E)=Hom^{0}_{H^{*}(B)-Mod^{gr}}(E,E[n])=0,$$
unless $E$ would be isomorphic to $E[n]$ which would a condradiction with the 
finite dimensionality of $E$. Let now $E' \in D(B)$ be an object such that 
$\phi(E')\simeq E$. As the graded functor $\phi$ is an equivalence, we have
$$[E',E']\simeq k \qquad [E',E'[n]]=0 \; if \, n\neq 0.$$
This implies that, the $k$-dg-algebra
$\mathbb{R}\underline{Hom}_{k}(E',E')$ is quasi-isomorphic to $k$, and thus that the pair of adjoint 
functors
$$E \otimes_{k} - : D(k) \longrightarrow D(B) \qquad 
\mathbb{R}\underline{Hom}_{k}(E,-) : D(B) \longrightarrow D(k)$$
defines a full embedding $D(k) \hookrightarrow D(B)$. Therefore, when considered
as a morphism $k \rightarrow B$ in $Hom_{k}$, the object $E$ exhibits $k$ as a retract of $B$. 
By dualizing $B$ we get that $B^{op}$ is a retract of $k$ in $Hom_{k}$. The object 
$B^{op}$ is then the image of a projector $p$ on $k$ in $Hom_{k}$, which corresponds to 
a compact object $P \in D(k)$ with the property that $P\otimes_{k}P \simeq P$. The only possibility is
$P\simeq k$, giving that the projector $p$ is in fact the identity. This implies that $B^{op}$ is isomorphic
to $k$ in $Hom_{k}$, and by duality that $B$ is isomorphic to $k$.  \hfill $\Box$ \\
\end{document}

% % \bigpablo{We can define refinements as: embedding, open embedding or isomorphism of topological spaces. The first is better to define algebraic data, the last better to define $(\infty,n)$-categories with adjoints.}

% % Recall the definitions of (conically smooth) stratified spaces.

% % \begin{definition}[Refinement]
% %     Let $f\in (X \to P) \to (Y \to Q)\in \Strat$.
% %     We say that $f$ is an \textbf{open refinement} if it is an open embedding of underlying topological spaces, and for each $p  \inset P$ the restriction $f_{|p}: X_p \to Y_{f(p)}$ is an embedded submanifold.
% % \end{definition}

% % \begin{definition} We define the ordinary $1$-category $\Ran_{\mathbb{R}^n}$ whose
% %     \begin{enumerate}
% %         \item \textbf{objects} are open refinements $X \to K \times \mathbb{R}^n$ such that $X \to K$ is a constructible bundle whose fibers are finitary,
% %         \item \textbf{morphisms} are stratified maps $X \to X'$ and $K \to K'$ inducing a pullback squares
% %         % https://q.uiver.app/#q=WzAsNCxbMCwwLCJYIl0sWzEsMCwiWCciXSxbMCwxLCJLXFx0aW1lcyBcXG1hdGhiYntSfV5uIl0sWzEsMSwiSydcXHRpbWVzIFxcbWF0aGJie1J9Xm4iXSxbMiwzXSxbMCwyXSxbMSwzXSxbMCwxXSxbMCwzLCIiLDEseyJzdHlsZSI6eyJuYW1lIjoiY29ybmVyIn19XV0=
% %         \[\begin{tikzcd}
% %             X & {X'} \\
% %             {K\times \mathbb{R}^n} & {K'\times \mathbb{R}^n}
% %             \arrow[from=1-1, to=1-2]
% %             \arrow[from=1-1, to=2-1]
% %             \arrow["\lrcorner"{anchor=center, pos=0.125}, draw=none, from=1-1, to=2-2]
% %             \arrow[from=1-2, to=2-2]
% %             \arrow[from=2-1, to=2-2]
% %         \end{tikzcd}\]
% %     \end{enumerate}
% %     We have a forgetful functor $\ran_{\mathbb{R}^n} \to \strat$.
% % \end{definition}

% % \begin{proposition}
% %     The forgetful $\ran_{\mathbb{R}^n} \to \strat$ is a transversality sheaf.
% %     We denote by $\Ran_{\mathbb{R}^n}$ the associated $(\infty,1)$-category.
% % \end{proposition}

% % We define a sub-category $\Bsc_{\mathbb{R}^n} \mono \Disk_{\mathbb{R}^n} \mono \Snglr_{\mathbb{R}^n} \mono \Ran_{\mathbb{R}^n}$ on maps being open-embeddings (to add details).

% % For $\B \to \Bsc_{/\mathbb{R}^n}$ a morphism of fibrations over $\Bsc$, we have $\textnormal{\textbf{Free}}_{\B} \in \EAlg{n}(\Cat{1})$ defined by the pullback
% % % https://q.uiver.app/#q=WzAsNCxbMCwwLCJcXHRleHRub3JtYWx7XFx0ZXh0YmZ7RnJlZX19X3tcXEJ9Il0sWzEsMCwiXFxEaXNrKFxcQnNjKV97L1xcbWF0aGJie1J9Xm59Il0sWzAsMSwiXFxEaXNrKFxcQikiXSxbMSwxLCJcXERpc2soXFxCc2Nfey9cXG1hdGhiYntSfV5ufSkiXSxbMiwzXSxbMSwzXSxbMCwxXSxbMCwyXSxbMCwzLCIiLDEseyJzdHlsZSI6eyJuYW1lIjoiY29ybmVyIn19XV0=
% % \[\begin{tikzcd}
% % 	{\textnormal{\textbf{Free}}_{\B}} & {\Disk(\Bsc)_{/\mathbb{R}^n}} \\
% % 	{\Disk(\B)} & {\Disk(\Bsc_{/\mathbb{R}^n})}
% % 	\arrow[from=1-1, to=1-2]
% % 	\arrow[from=1-1, to=2-1]
% % 	\arrow["\lrcorner"{anchor=center, pos=0.125}, draw=none, from=1-1, to=2-2]
% % 	\arrow[from=1-2, to=2-2]
% % 	\arrow[from=2-1, to=2-2]
% % \end{tikzcd}\]
% % And we also have $\textnormal{\textbf{Free}}^{\sifted}_{\B}\in \EAlg{n}(\Cat{1}^{\sifted})$.

% % We now define what ought to be thought about as the category of ``free $(\infty,n)$-categories with adjoints''.

% % \begin{definition}
% %     Consider $\Past_n$ to be the $(\infty,1)$-category whose:
% %     \begin{enumerate}
% %         \item \textbf{objects} are basics $\B \to \Bsc_{/\mathbb{R}^n}$
% %         \item \textbf{morphisms} from $\A \to \B$ are the data of:
% %         for all $a \in \A$,
% %         an object $X_a \in \Snglr_{\reln}$ together with a lift $\tau_a \to \B \to \Bsc_{\reln}$,
% %         a map $X_a \to U_a \in \Bun_{\reln}$.
% %     \end{enumerate}
% % \end{definition}

% % We have a map $\Theta_n \to \Past_n$ given by subdividing grids. And a map $\Past_n \to \EAlg{n}(\Cat{1}^{\sifted})$ given by $\textnormal{\textbf{Free}}^{\sifted}_{-}$

% % \begin{theorem}
% %     Every representable presheaf is a flagged $(\infty,n)$-category.
% % \end{theorem}

% \subsection{Codimension One Tangles}

% We use the definitions of higher bordism categories from \cite{CS19,Sch14}.

% \begin{definition}[{\cite[Section 4]{CS19}}] We define $\Int_k \in (\Delta^\op) \to \Top$ by
%     \[\begin{cases}[rCl]
%         (\Delta^\op) &\longto& \Top \\
%         {[k]} &\longmapsto& \Int_k \coloneqq \{(a,b)\inset (\mathbb{R}^{k+1})^2 \mid \forall 0 \leq i \leq k, a_i < b_i \textnormal{ and }(\forall 0 \leq i \leq k-1, a_i \leq a_{i+1} \textnormal{ and }b_i \leq b_{i+1})\}\\
%         (\iota\in [i] \to [j])^\op &\longmapsto& \begin{cases}[rCl]
%             \Int_j &\to& \Int_i\\
%             ((a_s)_{s=0}^j,(b_s)_{s=0}^j) &\mapsto& ((a_{\iota(s)})_{s=0}^i,(b_{\iota(s)})_{s=0}^i)
%         \end{cases}
%     \end{cases}\]
%     where the subsets are endowed with subspace topologies.
% \end{definition}

% We now slightly modify the definiton of $\PBord$ from \cite{CS19} to define (what ought to be) the ``free $(\infty,n)$-category with adjoints generated by a single $1$-cell''.

% \begin{definition} Let $([k_i])_{i = 1}^n \in (\Delta^\op)^n$. We define $\PAdj_{k_1, \ldots, k_n}^\infty$ to be the topological space whose points are the data of $(I, M, \rho, l)$ where:
% \begin{enumerate}
%     \item $I = ((a^i_s)_{s=0}^{k_i},(b^i_s)_{s=0}^{k_i})_{i=1}^n \inset \prod_{i=1}^n \Int_{k_i}$,
%     \item $M \mono \prod_{i=1}^n (a_0, b_{k_i})$ is an embedded submanifold of codimension $1$ such that: for every subset $S \subseteq \{1, \ldots, n\}$ and every $s \inset \prod_{i \inset S} [k_i]$, the projection $M\cap\prod_{i \inset S} (a_{s_i}, b_{s_i}) \to \prod_{i \inset S} (a_{s_i}, b_{s_i})$ is fiber bundle.
%     \item $\rho\in TM \oplus \underline{\mathbb{R}} \simeq \underline{\mathbb{R}}^n$ an $\mathbb{R}$-trivilization of the tangent bundle,
%     \item $l \in \pi_0(\prod_{i=1}^n (a_0, b_{k_i}) \setminus M) \to \{\pm\}$
% \end{enumerate}
% endowed with the finest topology making the following maps continuous\dots
% This data arranges into a $n$-fold simplicial space
% \[\begin{cases}[rCl]
%     (\Delta^\op)^n &\longto& \Spa \\
%     ([k_i])_{i = 1}^n &\longmapsto& \PAdj_{k_1, \ldots, k_n}^\infty
% \end{cases}\]
% \end{definition}

% \begin{proposition}
%     The $n$-fold simplicial space $\PAdj_{-,\ldots,-}^\infty$ satisfies the Segal conditions.
% \end{proposition}

% We denote by $\Adj_{-,\ldots,-}^\infty$ its univalent completion.

% \subsection{Sheaf Definition}

% \begin{definition}
%     Let $([k_i])_{i = 1}^n \in (\Delta^\op)^n$ and $X$ be a finitary manifold.
%     We define $\PAdj_{k_1, \ldots, k_n}^\infty(X)$ to be the set of $(M, \rho)$ where:
%     \begin{enumerate}
%         \item $I = ((a^i_s)_{s=0}^{k_i},(b^i_s)_{s=0}^{k_i})_{i=1}^n \inset \prod_{i=1}^n\Int_{k_i}(X)$.
%         \item $M \subseteq X \times \prod_{i=1}^n (a_0, b_{k_i})$ is an embedded submanifold of codimension $1$,
%         \item $\rho\in M^c \to \{\pm\}$ a continuous map,
%     \end{enumerate}
%     such that:
%     \begin{enumerate}
%         \item $M \to X$ is a fiber bundle,
%         \item for every subset $S \subseteq \{1, \ldots, n\}$ and every $s \inset \prod_{i \inset S} [k_i]$, the projection $M\cap\prod_{i \inset S} (a_{s_i}, b_{s_i}) \to \prod_{i \inset S} (a_{s_i}, b_{s_i})$ is fiber bundle.
%         \item for every $x \inset M$, we have an open set $x \inset U$ such that (separate between $\pm$)
%     \end{enumerate}
% \end{definition}

% \subsection{Truncated Definition}

% \begin{definition}
%     Let $([k_i])_{i = 1}^n \in (\Delta^\op)^n$.
%     We define $\PAdj_{k_1, \ldots, k_n}^\infty(X)$ to be the quotient set whose elements are $(M, \rho)$ where:
%     \begin{enumerate}
%         \item $M \subseteq \mathbb{R}^n$ is an embedded submanifold of codimension $1$,
%         \item $\rho\in \mathbb{R}^n\setminus M \to \{\pm\}$ a continuous map,
%     \end{enumerate}
%     such that:
%     \begin{enumerate}
%         \item for every subset $S \subseteq \{1, \ldots, n\}$ and every $s \inset \prod_{i \inset S} [k_i]$, the projection $M$ is transverse to $\prod_{i \inset S} {s}$,
%         \item for every $x \inset M$, we have an open set $x \inset U$ such that (separate between $\pm$)
%     \end{enumerate}
% \end{definition}

% \subsection{Excision}

% \begin{definition}
%     Write $S^0 = \{\pm\}$. Let $\Delta_{\leq 1}^\inj \to \Snglr$ be the diagram
%     % https://q.uiver.app/#q=WzAsMixbMCwwLCJcXG1hdGhiYntSfSJdLFsxLDAsIkMoU14wKSJdLFswLDEsIkMoLSkiLDIseyJjdXJ2ZSI6Miwic3R5bGUiOnsidGFpbCI6eyJuYW1lIjoiaG9vayIsInNpZGUiOiJib3R0b20ifX19XSxbMCwxLCJDKCspIiwwLHsiY3VydmUiOi0yfV1d
% \[\begin{tikzcd}
% 	{\mathbb{R}} & {C(S^0)}
% 	\arrow["{C(-)}"', curve={height=12pt}, hook', from=1-1, to=1-2]
% 	\arrow["{C(+)}", curve={height=-12pt}, from=1-1, to=1-2]
% \end{tikzcd}\]
%     Define a functor $\Delta^\inj \to \Snglr$  by left Kan extension:
%     in other words, it maps $[k]$ to the stratified space whose underlying topological space is $\mathbb{R}$ with $k$ marked points over the canonical connected strata poset. 
% \end{definition}

% For $(M, \rho)$ and $X$, we want a map
% \[
%     X \times \Disk_{/I} \to \Mfld_{n}
% \]

% \subsection{Higher Morita Category}

% \subsubsection{Scheimbauer}

% We follow the definitions from \cite{Sch14}.

% \begin{definition}
%     A \textbf{relative category} $(\cat{C}, \cat{W})$ is the data of a $1$-category $\cat{C}$ together with a saturated class of arrows $\cat{W} \mono \cat{C}$.
%     The associated $(\infty,1)$-category is the localization $\cat{C}[\cat{W}^{-1}]$.
% \end{definition}

% \begin{proposition}
%     Any $(\infty,1)$-category can be obtained from a relative category.
% \end{proposition}
% \pablo{Ideally I want a symmetric monoidal version of this}

% \begin{definition}
%     A \textbf{symmetric monoïdal} relative category is the data of a relative category $(\cat{C}, \cat{W})$ together with a symmetric monoidal structure on $\cat{C}$ which preserves class $\cat{W}$ on each variable.
% \end{definition}

% I need some kind of result saying: algebra for an operad can be obtained as strict algebras.

% This could be true only for model categories, that'd be okay.

% \begin{definition}
%     Let $(\cat{C}, \cat{W})$ be a symmetric monoidal relative category.
%     Let $([k_i])_{i = 1}^n \in (\Delta^\op)^n$.
%     We define $\Mor_{k_1, \ldots, k_n}(\cat{C}, \cat{W})$ to be the space of locally constant
% \end{definition}


% \subsubsection{Mine}

% We pick a standard startificaton of $(0, k+1)$ given by $1,\ldots, k$.
% For $\alpha\in [i] \to [j]$ we have
% \[
%     s < x < s+1 \longmapsto \alpha(s) + (x-s)(\alpha(s+1) - \alpha(s))
% \]
% For any length $l$ composition $\alpha_1 \circ \ldots \circ \alpha_l$, we find isotopies.

% This gives a $\infty$-functor $\Delta^\op \to $ the category of refinements.

% \begin{definition}
%     Let $([k_i])_{i = 1}^n \in (\Delta^\op)^n$.
%     We define $\Mor_{k_1, \ldots, k_n}^\infty(X)$ 
%     \begin{enumerate}
%         \item $((a_s^i)_{s=0}^{k_i})_{i=0}^n \in X \to \mathbb{R}$ with $a_s^i < a_{s+1}^i$,
%         \item $\rho\in \mathbb{R}^n\setminus M \to \{\pm\}$ a continuous map,
%     \end{enumerate}
%     such that:
%     \begin{enumerate}
%         \item for every subset $S \subseteq \{1, \ldots, n\}$ and every $s \inset \prod_{i \inset S} [k_i]$, the projection $M$ is transverse to $\prod_{i \inset S} {s}$,
%         \item for every $x \inset M$, we have an open set $x \inset U$ such that (separate between $\pm$)
%     \end{enumerate}
% \end{definition}

% We work with Claudia Scheimbauer's definition of the bordism category, processed through the $(0,1) \simeq \mathbb{R}$ diffeomorphism.
% A $k$-cell in $\Bord_n$ is an element of $\pi_0((\Bord_n)_{1,\ldots, 1, 0, \ldots, 0})$ with the first $k$-entries of the subscript being $1$.
% Note that by definition, we have a morphism of $n$-fold segal spaces $\PBord_n \to \Bord_n$  and hence a function $\pi_0((\PBord_n)_{1,\ldots, 1, 0, \ldots, 0}) \to \pi_0((\PBord_n)_{1,\ldots, 1, 0, \ldots, 0})$ (note that this map is not surjective).

% \begin{definition} A pre-$k$-cell in $\PBord_n$ is given by the choice of:
%     \begin{enumerate}
%         \item a finite dimensional vector subspace $V \subseteq \mathbb{R}^\infty$,
%         \item a $k$-manifold $M$,
%         \item an embedding $\iota \in M \mono V \times \mathbb{R}^k$ such that
%         \begin{enumerate}
%             \item $\pi \colon M \mono V \times \mathbb{R}^k \epi \mathbb{R}^k$ is proper,
%             \item for all $1 \leq i \leq k$, and for all $x \inset M$ such that $|\pi(x)_i| > 1$, $(\pi_i, \ldots \pi_n)$ is submersive at $x$. 
%         \end{enumerate}
%     \end{enumerate}
%     A sufficient condition for two pre-$k$-cells $(V_0, M_0, \iota_0)$ and $(V_1, M_1, \iota_1)$ to be equivalent is to have the data of:
%     \begin{enumerate}
%         \item a finite dimensional vector subspace $V \subseteq \mathbb{R}^\infty$ containing $V_1, V_2 \subseteq V$,
%         \item a $n$-manifold diffeomorphism $\phi\in M_0 \simeq M_1$,
%         \item a homotopy $\iota_{-}\in M_0 \times [0,1] \to V \times \mathbb{R}^k$ between $\iota_0$ and $\iota_1 \circ \phi$ (seen as maps to $V \times \mathbb{R}^k$) such that: for all $0 < t < 1$, the map $\iota_t \in M_1 \mono V \times \mathbb{R}^k$ is an embedding satisfying
%         \begin{enumerate}
%             \item $\pi \colon M_0 \mono V \times \mathbb{R}^k \epi \mathbb{R}^k$ is proper,
%             \item for all $1 \leq i \leq k$, and for all $x \inset M$ such that $|\pi(x)_i| > 1$, $(\pi_i, \ldots \pi_k)$ is submersive at $x$. 
%         \end{enumerate}
%     \end{enumerate}
% \end{definition}

% \begin{lemma}
%     Any pre-$k$-cell $(V, M, \iota)$ gives a point in $[(V\times \mathbb{R}^{n-k}, M\times \mathbb{R}^{n-k}, \iota \times \mathbbm{1}_{\mathbb{R}^{n-k}}, (-\infty, -1) \leq (1, \infty))] \inset \pi_0((\PBord_n)_{1,\ldots, 1, 0, \ldots, 0})$ and hence in $\pi_0((\Bord_n)_{1,\ldots, 1, 0, \ldots, 0})$.
% \end{lemma}

% \begin{definition}[$k$-handles]
%     Let $1 \leq k \leq n$ and $\sigma \inset \{\epsilon, \eta\}^k$ be a sequence of letters.
%     We define a pre-$k$-cell $\sigma(+)$ of $\Bord_n$ by $(\mathbb{R}^k, \mathbb{R}, \iota_\sigma)$ where
%     \[\iota_\sigma \in \begin{cases}[rCl]
%         \mathbb{R}^k & \longmono & \mathbb{R} \times \mathbb{R}^k \\
%         (x_1, \ldots x_k) & \longmapsto & \left(x_1, \ldots x_k, \sum_{i=1}^k |\sigma_i| x_i^2 \right)
%     \end{cases}\]
%     where $|\epsilon| = -1$ and $|\eta| = 1$, with the canonical framing.
% \end{definition}

% \begin{lemma}[Handlebody Decomposition]
%     Let $M$ be a pre-$k$-cell in $\PBord_n$.
%     Suppose that the projection $\pi_k$ is a Morse function.
%     Then one can see $M$ as a $k$-composition of handles. 
% \end{lemma}
% \begin{proof}
%     Let $(x_k)_1, \ldots (x_k)_l$ be the critical values.
%     By the Morse lemma, we have coordinate systems 
% \end{proof}